%!TEX root = ../../main.tex

\chapter{Einleitung}
\label{ch:einleitung}

\section{Motivation und Problemstellung}
\label{sec:motivation}
% Beschreibe hier den Kontext: Bedeutung von Daten in Unternehmen,
% Herausforderungen klassischer ETL-Prozesse, Aufkommen von LLMs.

\section{Zielsetzung und Forschungsfragen}
\label{sec:zielsetzung}
% Formuliere das Ziel der Arbeit und 2-4 konkrete Forschungsfragen.
% Beispiel:
% \begin{itemize}
%     \item FF1: Inwiefern können LLMs klassische ETL-Methoden ergänzen oder ersetzen?
%     \item FF2: Welchen Einfluss haben unterschiedliche Prompting-Strategien auf die Ergebnisqualität?
%     \item FF3: Welche Grenzen zeigen sich beim Einsatz von LLMs in ETL-Prozessen?
% \end{itemize}

\section{Methodisches Vorgehen}
\label{sec:methodik}
% Kurzer Überblick über das experimentelle Vorgehen.

\section{Aufbau der Arbeit}
\label{sec:aufbau}
% Kurze Vorstellung der nachfolgenden Kapitel.
